%! TeX root = ../main.tex

\sectionanchor{toc:biography}{A Life Remembered}

% ---------------------------------------------------------------------------
% Two pages, on the pattern the Photographs section already sets: a full-bleed
% frontispiece, then the text.
%
% The frontispiece is plates/front-life.png -- the agbada portrait, the
% youngest picture of him the family has, cover-cropped to the page with a
% white field burnt into the foot for the type. It is built by
% assets/make-plates.sh, from plate-agbada, by the same `front' function that
% builds the Photographs frontispiece; to change the picture, change which
% plate that function is handed there and re-run the script.
%
% A different photograph from the one that opens Photographs, deliberately.
% Two sections opening on the same portrait reads as a shortage of pictures
% rather than as a motif -- and the young man in the agbada is the right one
% to stand at the head of a page that begins with his birth.
% ---------------------------------------------------------------------------

\pageghost{front-life}

\vspace*{\fill}

\begin{center}
  {\Large\scshape\color{primary} A Life Remembered}

  \vspace{3mm}
  {\color{rulecolor}\rule{18mm}{0.4pt}}

  \vspace{4mm}
  \begin{minipage}{0.84\linewidth}
    \centering\small\itshape\color{accent}
    ``Onoteagbor'' --- the one who would live long
  \end{minipage}
\end{center}

\vspace*{6mm}

\newpage

% ---------------------------------------------------------------------------
% The biography itself. The ledger ghost behind it, which is the photograph of
% him taking stock: he was, by his family's account, an exact keeper of
% records, and it is the one wash in the set that says so.
%
% Their own first line -- "BIOGRAPHY OF THE LATE SIR FRANCIS ONOTEAGBOR
% ATSEKHAMEH (FRANCO NERO THE SKY MAN)" -- is the title of the piece rather
% than a sentence of it, so it is set as one: the name in letterspaced caps
% and the two names he was known by beneath it.
% ---------------------------------------------------------------------------

\pageghost{ghost-ledger-hi}

\begin{center}
  {\coverface\fontsize{11}{14}\selectfont\color{primary}
  \coverline[150]{Sir Francis Onoteagbor Atsekhameh}\par}

  \vspace{1.5mm}
  {\small\itshape\color{accent} ``Franco Nero'' \diamondmark\ De Sky Man}
  \thinseparator
\end{center}

\vspace{2mm}

% Justified, with an indent, and not the ragged setting the liturgy pages
% take. The liturgy is a list of short spoken lines and rags cleanly; this is
% continuous prose in a 57mm column, where ragged right breaks into stubs.
% The indent rather than a blank line between paragraphs for the same reason
% the tributes use one: at this length the paragraph gaps would cost a
% centimetre of page and buy nothing an indent does not.
\begin{multicols}{2}
  \small
  \setlength{\parindent}{1.2em}
  \setlength{\parskip}{0.5pt}
  \raggedcolumns
  \color{primary}
  % A 57mm column at 10pt holds about seven words, and TeX's default
  % tolerance gives up on a line and stretches the spaces before it has tried
  % everything. Raising it, and leaving a little emergency stretch for the
  % paragraph it still cannot set, buys back most of the loose lines. Set
  % here rather than in the preamble: the liturgy pages are ragged right and
  % would not notice, and nothing else in the booklet should reflow for this.
  \tolerance=1600
  \emergencystretch=1.2em

  \lifehead{His Early Years}

  \dropcap{T}{he late} Francis Onoteagbor \\Atsekhameh was born into the family
  of Otse to Late Mr.\ Apuede and Mrs.\ Egieobo Atsekhameh in the year 1946.

  He was the first surviving child of his parents amongst 7, hence the name
  \textsc{Onoteagbor} which translates to --- \textit{The One Who Would Live
  Long}.\smallskip

  As a little boy, who lost his father at a very tender age, Franco
  Nero refused to accept that he was born by his earthly parents, so he came
  up with a narrative saying \textit{``he fell from the sky''}, and his mother was only
  privileged to catch him with a ``basin'' and he proceeded to call himself
  ``\textsc{Adeokwile}'' which literally translates to someone who fell from
  the sky. He took this narrative very seriously, told everyone in his family
  and the entire community and insisted to be addressed by that.

  \lifeline{Hence the name ``\coverline[150]{Sky Man}''.}

  \lifehead{His Education}

  Franco Nero began his elementary education at St.\ Thomas Aquinas Catholic
  Primary School, Ivioghe, Agenebode. After completing his primary education,
  he went on to live with his older cousin, High Chief Thomas Audu, in Sapele,
  Delta State. Who later relocated to Lagos, where he went on to further his
  education through distant learning with the prestigious London Business
  School in England, United Kingdom, and obtained a Diploma in Sales
  Management. He did't just stop there; he pushed himself hard to further obtain an Advanced Diploma in Marketing.

  \lifehead{His Career}

  He began his career with the Ministry of Defence and then proceeded to work
  as a Regional Sales Manager with Omon Stores, and later got employed by the
  then famous British Trading Company that played a massive role in West
  African commerce, particularly in Nigeria and Ghana --- G.B Ollivant (GBO). GBO was ultimately absorbed by the United Africa Company (UAC), and later
  became a part of the Unilever Group --- De \textsc{Sky Man} capitalized on
  these corporate growth opportunities, leveraging the organizational
  support to expand his professional skillsets, and made the sky his
  springboard, because he excelled greatly in his Salesmanship career.

  He also worked with Pilroads Nigeria Ltd before retiring and venturing into
  several businesses. He was a serial businessman who explored so many
  business ideas in his lifetime.

  Despite losing his father at a very tender age, he didn't let his
  circumstances hinder the man he was destined to be. His determination and
  drive to become a person of great value, was reflected throughout his
  lifetime.

  \lifehead{His Hobbies}

  Sky Man loved singing, dancing, listening to music, reading novels, watching the news, and storytelling.

  He was an astute lover of sports, particularly football and wrestling,

  and he had a deep passion for Arsenal, a club he proudly cheered on for well
  over 3 decades.

  Our Dad was a very meticulous and detailed-oriented man, who was very
  particular about his environment. He was also an excellent keeper of
  records, and was a champion at recounting events with dates. He penned down
  the birthdays of all his children and grandchildren, including important
  anniversary dates. He was a very smart and intelligent man with a stellar
  memory.

  \lifehead{His Faith}

  As a Christian and a steadfast catholic, his love for Mother Mary was
  admirable. He was a very devoted legionary who headed Praesidiums and
  carried out his legion work in the remotest villages within and outside
  Benin City. He would evangelise to unbelievers and idol worshipers, tell
  them about Jesus Christ, and fearlessly destroyed several shrines in a bid
  to convert the owners and lead them to Christ.

  It was no coincidence he died on August 15\textsuperscript{th} --- the day
  the church celebrates the feast of Assumption.

  \lifehead{His Personality}

  Daddy had quite a large and accommodating heart. He opened the door of our
  house to everyone --- his family, his wife's relatives, his children's friends,
  and pretty much all who needed a home at different times, and he never complained of
  anyone's presence or how long they lived with us.

  He loved action movies, loved reading and writing, dancing and everything
  that had to do with music --- listening, singing, and dancing, and oh, he
  had a great taste in music. Dad didn't joke with his oldies, we still recall
  how he would write the complete lyrics of songs on papers for us to read and
  sing along with him. He was blessed with a sweet melodious voice, that
  charged the atmosphere in a very awesome way, whenever it rents the air.

  Daddy loved fashion till his last days. He enjoyed rocking his sneakers,
  jeans, t-shirts, and his face caps, and would always joke about being a
  ``small boy'' when asked.

  \lifehead{His Marriage and Family}

  Sky Man got married to his ``beautiful wife'' and friend as he liked to call
  her on the 22\textsuperscript{nd} day of May, 1977, and their union was
  blessed with 5 children.

\end{multicols}

% ---------------------------------------------------------------------------
% From here the family stop describing him and speak to him, and the setting
% changes with the voice: out of the columns, full measure, in italic. See
% \lifeenvoi in main.tex.
%
% \newpage rather than letting it fall where it falls. The columns end a few
% lines short of the foot -- not enough for the envoi, which then broke to
% the next page and left the \vspace*{\fill} meant to centre it behind on
% the page before, absorbed into the gap at the bottom of the second column.
% Asking for the page explicitly is what makes the two fills here a pair, and
% what centres the family's last words on a page of their own.
% ---------------------------------------------------------------------------

\newpage

% The photograph behind the farewell: him saluting, in the Legion's green,
% under the light. It runs to the trim on all four sides -- \pageghost is the
% same macro that places the cover and the frontispieces -- and it is veiled
% so the words can be set straight onto it rather than into a cleared field.
%
% The veil is baked into the plate by assets/make-plates.sh, not applied
% here, and which photograph it uses is set in assets/data/plates.toml under
% [close-life], with the strength of the veil and where it reaches full
% strength. Change the picture there and re-run the script; this line does
% not move.
\pageghost{close-life}

\vspace*{\fill}

\begin{lifeenvoi}
  Today, our hearts are broken as we speak about you in past tenses, though
  you are not physically present, but the memories we created and shared will
  forever remain precious to our hearts.

  We believe you are in a better place, with your father in heaven, we know
  you are singing, dancing and definitely cracking up the hosts of heavens.
  Most importantly, we know you're watching over us, and we are consoled that
  you get to be with the one who loves you most.

  Daddy, thank you for the unconditional love, the sacrifices, the priceless
  moments, and the beautiful memories created and shared. They will live in
  our hearts and be cherished forever.
\end{lifeenvoi}

\lifeclose{Continue to rest in the bosom of the Almighty,}%
{Franco Nero, De Sky Man}

\vspace*{\fill}
